% =====================================================================
% Paper B: structure only (headings, planned figures and tables).
% Generated 2026-10-09 from main.tex and sections/03_data_protocol.tex (commit 0fc16c7).
% Keep in sync by hand when a heading changes. Build: pdflatex structure
% =====================================================================
\documentclass[9pt,journal,a4paper]{IEEEtran}

\usepackage{xcolor}
\usepackage[hidelinks]{hyperref}

% Planned figure or table, shown in place.
\newcommand{\planned}[2]{\noindent\textcolor{gray}{\footnotesize\textit{[#1] #2}}\par}

\begin{document}

\title{Data-Driven Quantification of Aggregated Thermostatically Controlled Loads from Low-Voltage Net Load: What Can Be Learned?}
\author{Alexandre~M.~V.~Gouveia, Md. Umar Hashmi, Reinhilde D'hulst, Dirk Van Hertem}
\maketitle

\begin{abstract}
\textcolor{gray}{\footnotesize\textit{Written last.}}
\end{abstract}

\section*{Nomenclature}

\section{Introduction}
\planned{Table}{Literature: method family, input, output, labels needed, capacity, transfer tested}

\section{Problem Statement}
\subsection{Aggregate Model}
\subsection{Quantities of Interest}
\planned{Table}{Targets: definition, unit, DSO use, determined by the net load}
\subsection{What the Net Load Determines}
\subsection{Information Available to a DSO}

\section{Data and Evaluation Protocol}
\subsection{Source Datasets}
\planned{Table}{Source datasets (HEAPO, Kaiser et al., EoH, LCL)}
\subsection{CH and GB Datasets}
\planned{Table}{The CH and GB datasets after selection}
\subsubsection{CH dataset}
\subsubsection{GB dataset}
\subsubsection{Labels per dataset}
\subsection{Semi-Synthetic Aggregates}
\planned{Fig.}{Construction of the semi-synthetic aggregates and household-disjoint splits}
\subsubsection{Fillers for the GB dataset from analog days}
\subsection{Evaluation Protocol}
\subsection{Pitfalls of Semi-Synthetic Benchmarks}
\planned{Table}{Legacy protocol vs leak-free protocol}

\section{Estimators}
\subsection{Physics-Based Estimators}
\subsection{Anchor-Only Baselines}
\subsection{Features}
\subsection{Residual Learning on the Physics Estimate}
\subsection{Learning Models}
\subsection{Learning the Transfer Parameter}

\section{Targets and Identifiability (RQ1)}
\subsection{Accuracy per Target}
\planned{Fig.}{Error per target, with and without scale anchor}
\subsection{Role of the Scale Anchor}
\subsection{Large-Sample Check at Daily Resolution}

\section{Value over the Physics Estimator (RQ2)}
\subsection{Benchmark on Semi-Synthetic Aggregates}
\planned{Table}{Method $\times$ anchor, median and 5--90\% band across splits}
\subsection{Gain over Anchor-Only Baselines}
\subsection{Dependence on Penetration and Aggregate Size}
\planned{Fig.}{Error vs HP penetration and aggregate size (ML, physics, anchor-only)}
\subsection{Does More Training Data Help?}
\subsection{Error Decomposition}
\subsection{Other Electric Heating and Utility Control}

\section{Transfer across Populations (RQ3)}
\subsection{Transfer Settings}
\subsection{End-to-End Models}
\subsection{Learned Transfer Parameter}
\planned{Table}{Transfer error per target population, SF transfer vs learned parameter}
\subsection{Recalibration with a Small Local Pilot}

\section{Detecting Change in TCL Capacity (RQ4)}
\subsection{Uptake Scenarios}
\subsection{Multi-Year Aggregates}
\subsection{Detection and Quantification}
\planned{Fig.}{Minimum detectable change vs aggregate size and horizon}
\subsection{Real Change on FeederBW}

\section{Discussion}
\subsection{What the Models Learn}
\subsection{Accuracy versus Data Cost}
\subsection{Guidance for DSOs}
\subsection{Limitations}

\section{Conclusion}

\appendices
\section{Household Datasets and Design Envelope}
\planned{Table}{Design envelope of the semi-synthetic aggregates}
\section{Model Hyperparameters}

\end{document}
